\documentclass[11pt]{article}
\usepackage[margin=1.1in]{geometry}
\usepackage{times}
\usepackage[T1]{fontenc}
\usepackage{setspace}
\onehalfspacing
\usepackage{parskip}
\pagenumbering{gobble}

\title{Earth Is a Protected Historical Site}
\author{Flyxion}
\date{September 2026}

\begin{document}
\maketitle

The campaign to save the planet succeeded, in the narrow technical sense the campaign had actually been measuring. Emissions fell. The designation went through. Earth is now listed, in the appropriate registries, as a protected heritage site of the first order, and the listing did exactly what listings do: it raised the cost of everything within the boundary, the land, the water rights, the permits required to keep a house standing on soil now classified as irreplaceable, until the protection accomplished what several decades of unprotected neglect had not, which was making the planet too expensive for most of the species it was protected for to go on living on.

Nobody planned this as the mechanism, which is what makes it worth noting rather than merely regrettable. The instruments were all individually reasonable: a carbon levy, a conservation easement, a heritage tax funding restoration work, each one justified on its own terms and each one, in combination, pricing the original inhabitants out of the thing being conserved on their behalf, the way a neighborhood's most passionate preservationists are so often the first to be unable to afford the neighborhood once the preservation succeeds. The remaining population, the ones who couldn't relocate or wouldn't, are absorbed into the heritage economy the designation created, employed to greet visitors, narrate the old way of living, demonstrate a traditional dinner or a traditional commute for tour groups arriving from stations in orbit, where the actual cost of living turned out to be lower than it is on the surface of the planet the tour is about.

The visitors leave moved. The planet is, by every metric the campaign originally cared about, saved. The people who used to simply live there now perform living there, for a wage, on a schedule, under lighting chosen for the tour, and if anyone still calls this home, it is mentioned in the brochure as one of the more affecting details of the authentic experience on offer.

\end{document}
